\subsubsection{Perfect fields}
\label{editorial-source-section-perfect-fields}

\noindent
Here is the definition.

\begin{definition}
\label{editorial-source-definition-perfect}
Let $k$ be a field. We say $k$ is {\it perfect}
if every field extension of $k$ is separable over $k$.
\end{definition}

\begin{lemma}
\label{editorial-source-lemma-perfect}
A field $k$ is perfect if and only if it is a field of characteristic $0$
or a field of characteristic $p > 0$ such that every element has a $p$th
root.
\end{lemma}

\begin{proof}
In characteristic zero every field extension is
separable, by the characteristic-zero argument in
Lemma \ref{editorial-source-lemma-separably-generated-separable}.
This is precisely the stated definition of perfect.

Suppose now that $\operatorname{char}(k)=p>0$.
If $k$ is perfect, take any $a\in k$ and its unique
$p$th root $b$ in an algebraic closure. The original
extension $k(b)/k$ is purely inseparable, since
$b^p=a$, and is separable by perfectness. A minimal
polynomial of $b$ divides $T^p-a=(T-b)^p$ over the
algebraic closure. Thus it has only one distinct
root, whereas separability makes all its roots
distinct. Its degree is therefore one, so $b\in k$.
This includes $a=0$, whose root is zero.

Conversely, suppose every $a\in k$ has a $p$th root
in $k$. Frobenius is injective on a field, so those
roots are unique and the specified extension
$k^{1/p}/k$ is the identity. For every original
field extension $K/k$ the tensor map
$$
K\otimes_k k\longrightarrow K,\qquad
\lambda\otimes a\longmapsto a\lambda
$$
has inverse $\lambda\mapsto\lambda\otimes1$.
It identifies $K\otimes_k k^{1/p}$ with the reduced
field $K$. Lemma
\ref{editorial-source-lemma-characterize-separable-field-extensions}
then makes every original $K/k$ separable, so $k$
is perfect.

In particular, in positive characteristic it
suffices to test the separability of the individual
extensions $k(a^{1/p})/k$. If $a$ has no root in $k$,
then $T^p-a$ is irreducible: a proper monic factor
over $k$ would be $(T-b)^d$ over the algebraic
closure for some $1\leq d<p$, whose coefficient
of $T^{d-1}$ is $-db$. Since $d\neq0$ in $k$,
that coefficient would put $b$ in $k$, a contradiction.
Thus every nontrivial one of these tests has degree
exactly $p$ and is inseparable. This proves the
test and its exact degree without assuming a
general extension is finitely generated.
\end{proof}

\begin{lemma}
\label{editorial-source-lemma-make-separable}
Let $K/k$ be a finitely generated field extension.
There exists a diagram
$$
\xymatrix{
K \ar[r] & K' \\
k \ar[u] \ar[r] & k' \ar[u]
}
$$
where $k'/k$, $K'/K$ are finite purely inseparable field
extensions such that $K'/k'$ is a separable field extension.
In this situation we can assume that $K' = k'K$ is the compositum,
and also that $K' = (k' \otimes_k K)_{red}$.
\end{lemma}

\begin{proof}
Lemma \ref{editorial-source-lemma-make-separably-generated} supplies
the original diagram, with its upper-right field
temporarily denoted by $K'_0$, in which $k'/k$
and $K'_0/K$ are finite purely inseparable and
$K'_0/k'$ is separably generated. By Lemma
\ref{editorial-source-lemma-separably-generated-separable}, the last
extension is separable. Retain both original
embeddings into $K'_0$ and let $L=k'K\subseteq K'_0$
be their compositum. Since $k'\subseteq L\subseteq K'_0$,
Lemma \ref{editorial-source-lemma-subextensions-are-separable} makes
$L/k'$ separable. Also $L/K$ is finite purely
inseparable, being an intermediate extension of
$K'_0/K$. Thus choosing the new upper-right field
to be $K'=L$, with its actual inclusion into
$K'_0$, proves the compositum assertion.

For the tensor assertion, first treat characteristic
zero. A purely inseparable extension is then the
identity, so $k'=k$ and $K'_0=K$. The multiplication
map $k\otimes_k K\to K$, $a\otimes b\mapsto ab$,
is an isomorphism with inverse $b\mapsto1\otimes b$.
Its nilradical is zero, and it is the required map.

In characteristic $p>0$, put $B=k'\otimes_k K$.
Because the original extension $k'/k$ is finite
purely inseparable, there is a single $q=p^e$
with $(k')^q\subseteq k$: choose finitely many
field generators, take the largest of their
inseparability exponents, and apply Frobenius
to their full rational expressions. Retain
the multiplication map
$$
\mu:B\longrightarrow K'_0,\qquad
\sum_i a_i\otimes b_i\longmapsto\sum_i a_i b_i.
$$
It is balanced because the original field diagram
commutes. For every original $z=\sum_i a_i\otimes b_i$
one has the full identity
$$
z^q=\sum_i a_i^q\otimes b_i^q
=1\otimes\left(\sum_i a_i^q b_i^q\right),
\qquad
\mu(z)^q=\sum_i a_i^q b_i^q\in K.
$$
The equality for $q$ follows by iterating the
$p$th-power binomial identity; each intermediate
mixed coefficient is divisible by $p$. The map
$K\to B$, $b\mapsto1\otimes b$, is injective.
If $\mu(z)=0$, the coefficient sum in $K$ is
zero, and hence $z^q=0$. Conversely every
nilpotent maps to zero in the field $K'_0$.
Therefore $\ker\mu=\operatorname{Nil}(B)$.

The image of $\mu$ is the subring of finite sums
of products $a_i b_i$. It is a field: for a
nonzero such sum $w$, its $q$th power is a
nonzero element of $K$, and
$w^{-1}=w^{q-1}/w^q$ belongs to the same subring.
The image contains $k'$ and $K$, so it is exactly
the original $L=k'K$. Consequently the original
map gives the canonical $K$-algebra isomorphism
$$
(k'\otimes_k K)_{red}\longrightarrow L,
\qquad [\sum_i a_i\otimes b_i]\longmapsto\sum_i a_i b_i.
$$
Its inverse sends a finite sum in $L$ to the
class of the corresponding tensor; if two such
sums are equal, their difference lies in the
kernel just calculated, proving independence
of the representation. Thus both asserted
choices of $K'$ agree through this exact map.

This last comparison needs neither finite
generation of $K/k$ nor finiteness of the purely
inseparable extension in the first tensor factor.
Indeed let $l/k$ be any algebraic purely
inseparable extension and embed it, by its unique
roots, in an algebraic closure of the original
$K$. Every individual tensor in $l\otimes_k K$
uses finitely many coefficients $a_i\in l$,
so there is a $q=p^e$ for that finite list with
all $a_i^q\in k$. The same two displayed power
identities prove that the multiplication kernel
is exactly the nilradical, and the same inverse
formula for each nonzero image proves that the
image is the original compositum $lK$.
Thus $(l\otimes_k K)_{red}\simeq lK$ with the
same explicit map in this full generality.
Only the exponent may now depend on the tensor.

Finite generation cannot, however, be discarded
from the earlier assertion that one finite
purely inseparable base extension makes the
whole receiving extension separable. Take
$$
k=\mathbf F_p(t),\qquad
K=\bigcup_{r\geq0}\mathbf F_p(t^{1/p^r})
$$
inside a fixed algebraic closure, with $t$
transcendental and all root inclusions retained.
Every finite purely inseparable $k'/k$ has a
unique image in $K$: a root of an original
rational function $A(t)/B(t)$ of order $p^r$
is $A(t^{1/p^r})/B(t^{1/p^r})$, since its
coefficients belong to $\mathbf F_p$.
The elements of $k'$ have
$p^e$th powers in $k$ for some finite $e$.
The element $t^{1/p^{e+1}}$ is not in $k'$:
otherwise its $p^e$th power $t^{1/p}$ would
belong to $k$. The latter is impossible since
a rational function in $t$ cannot have $p$th
power $t$. To see this directly, writing such
a function as $A(t)/B(t)$ would give
$A(t)^p=tB(t)^p$; the exponents on the left
are divisible by $p$ and those on the right
are congruent to $1$ modulo $p$, so a nonzero
polynomial equality is impossible.
Thus $K/k'$ is a nontrivial purely inseparable
extension and cannot be separable. Its
compositum with the original $k'$ is still
$K$. Moreover any purported larger receiving
field separable over $k'$ would have the
subextension $K/k'$ separable, again impossible.
This proves the necessity of some restriction
on the original $K/k$ for the finite-base-change
existence assertion, while retaining the
unrestricted tensor-quotient comparison.
\end{proof}

\begin{lemma}
\label{editorial-source-lemma-perfection}
\begin{slogan}
Every field has a unique perfect closure.
\end{slogan}
For every field $k$ there exists a purely inseparable extension
$k'/k$ such that $k'$ is perfect. The field extension
$k'/k$ is unique up to unique isomorphism.
\end{lemma}

\begin{proof}
In characteristic zero $k$ itself is perfect by
Lemma \ref{editorial-source-lemma-perfect}, and a purely inseparable
extension must be the identity. This gives
existence and uniqueness, including the specified
embedding of $k$.

Let $\operatorname{char}(k)=p>0$, and fix an
algebraic closure $\Omega$ of the original $k$.
For every integer $n\geq0$ put
$$
E_n=\{b\in\Omega:b^{p^n}\in k\}.
$$
The Frobenius identities show that this is a
field: sums and products have $p^n$th powers
in $k$, as does the inverse of a nonzero
element. One has $E_0=k$ and $E_n\subseteq E_{n+1}$.
Every $a\in k$ has a unique $p^n$th root in
$\Omega$, because $T^{p^n}-a$ has a root there
and the difference of any two roots has zero
$p^n$th power in a field. Thus $E_n$ has exactly
properties (a) and (b) of the stated construction.

To compare this with the original abstract
construction, let $k_n$ denote a copy of the
underlying field $k$, with structural map
$i_n:k\to k_n$, $a\mapsto a^{p^n}$.
The transition map $f_n:k_n\to k_{n+1}$ is
$r\mapsto r^p$, and
$f_n(i_n(a))=a^{p^{n+1}}=i_{n+1}(a)$.
The map
$$
\theta_n:k_n\longrightarrow E_n,\qquad
r\longmapsto r^{1/p^n}
$$
is a field isomorphism. Uniqueness of roots
proves additivity and multiplicativity;
surjectivity follows from the definition of
$E_n$, and injectivity follows from Frobenius.
It respects the original $k$-structure since
$\theta_n(i_n(a))=a$, and it respects the
transitions since
$\theta_{n+1}(r^p)=\theta_n(r)$.
Accordingly the union of the actual subfields
$E_n$ realizes the directed colimit of the
original copies with their Frobenius transition
maps. No literal inclusion of a copy with its
different structural map is being assumed.

Set $k'=\bigcup_{n\geq0}E_n$. This is a field,
and every one of its elements has a $p$-power
in the original $k$, so $k'/k$ is purely
inseparable. If $b\in E_n$, its unique $p$th
root in $\Omega$ lies in $E_{n+1}$, since its
$p^{n+1}$th power equals $b^{p^n}\in k$.
Thus Frobenius on $k'$ is surjective, and
Lemma \ref{editorial-source-lemma-perfect} makes $k'$ perfect.

For uniqueness, let $L/k$ be any perfect purely
inseparable extension, with its original
structural embedding of $k$. Every $a\in k$
has a unique $p^n$th root in $L$. Define
$\phi_n:E_n\to L$ by sending $b$ to that
root of $b^{p^n}\in k$. Frobenius and root
uniqueness prove that $\phi_n$ preserves
addition, multiplication, $1$, and the given
embedding of $k$. It is injective as a
nonzero field homomorphism. These maps agree
on the inclusions $E_n\subseteq E_{n+1}$,
again by root uniqueness, and give
$\phi:k'\to L$. Every $c\in L$ has some
$c^{p^n}\in k$, so it is the image of the
unique corresponding root in $E_n$.
Hence $\phi$ is surjective. Any $k$-homomorphism
from $k'$ to $L$ must take every such root
to that same unique root, so it equals $\phi$.
This proves uniqueness up to unique
$k$-isomorphism.

The construction proves the stronger universal
property with its maps: for every perfect
field $P$ and every field homomorphism
$\iota:k\to P$, there is exactly one field
homomorphism $\iota^{perf}:k'\to P$ extending
$\iota$. On $b\in E_n$ it is the unique
$p^n$th root in $P$ of $\iota(b^{p^n})$.
The preceding addition, multiplication and
compatibility proofs apply verbatim; they
did not use pure inseparability of the
receiving field for existence, only for
surjectivity. This also proves functoriality:
two composites that agree on $k$ send every
root to the same unique root. In
characteristic zero the assertion is just
the original map $\iota$ on $k'=k$.
\end{proof}

\begin{definition}
\label{editorial-source-definition-perfection}
Let $k$ be a field. The field extension $k'/k$ of Lemma \ref{editorial-source-lemma-perfection}
is called the {\it perfect closure} of $k$. Notation $k^{perf}/k$.
\end{definition}

\noindent
Let $l/k$ be any algebraic purely inseparable
extension. In positive characteristic, send an
original $a\in l$, with $a^{p^n}\in k$, to
the unique root of $T^{p^n}-a^{p^n}$ in
$k^{perf}$. The value is independent of the
chosen exponent, and a common exponent for
two elements proves additivity and
multiplicativity by Frobenius. This defines
the unique $k$-embedding $l\to k^{perf}$.
In characteristic zero $l=k$ and this is
the identity. Retain this actual embedding.
The extension $k^{perf}/l$ is purely
inseparable because every element has a
$p$-power in $k\subseteq l$, and its upper
field is perfect. It is therefore a
perfect closure of $l$. More explicitly,
the universal property just proved gives
maps $k^{perf}\to l^{perf}$ extending
$k\to l^{perf}$ and $l^{perf}\to k^{perf}$
extending the displayed embedding of $l$.
The first map fixes the original $l$
under these embeddings, since roots of
its elements' powers in $k$ are unique.
Both composites consequently fix the
base fields and every adjoining root,
so are identities. These are the
specified inverse isomorphisms of
perfect closures, not merely an
abstract equality of field types.

\begin{lemma}
\label{editorial-source-lemma-perfect-reduced}
Let $k$ be a perfect field.
Any reduced $k$-algebra is geometrically reduced over $k$.
Let $R$, $S$ be $k$-algebras.
Assume both $R$ and $S$ are reduced.
Then the $k$-algebra $R \otimes_k S$ is reduced.
\end{lemma}

\begin{proof}
By Lemma \ref{editorial-source-lemma-perfect}, either $k$ has
characteristic zero or its Frobenius map is
surjective. In characteristic zero the test
in Lemma
\ref{editorial-source-lemma-geometrically-reduced-finite-purely-inseparable-extension}
is precisely reducedness of the original $S$.
In positive characteristic $k^{1/p}=k$ through
the actual root embedding, and the same test
is reducedness of $k\otimes_k S\simeq S$,
with maps $a\otimes s\mapsto as$ and
$s\mapsto1\otimes s$. Thus every reduced
$k$-algebra is geometrically reduced. Applying
Lemma \ref{editorial-source-lemma-geometrically-reduced-any-reduced-base-change}
to the original reduced $R$ and geometrically
reduced $S$ proves that $R\otimes_k S$ is
reduced, including either zero-ring case.

The hypothesis on $k$ is also necessary for
this conclusion to hold for every pair of
reduced $k$-algebras, or even for every pair
of fields over $k$. If $k$ is not perfect,
Lemma \ref{editorial-source-lemma-perfect} gives characteristic
$p>0$ and an $a\in k$ with no $p$th root in
$k$. Retain $L=k(b)$, $b^p=a$. As proved in
Lemma \ref{editorial-source-lemma-perfect}, $T^p-a$ is
irreducible, so $1,b,\ldots,b^{p-1}$ is a
$k$-basis of $L$. The original map
$$
L[T]/(T^p-a)\longrightarrow L\otimes_k L,
\qquad T\longmapsto1\otimes b
$$
fixes $L$ through the left factor and is an
isomorphism by the corresponding bases.
Its inverse sends $c\otimes\sum_i d_i b^i$
to $\sum_i cd_i T^i$, for $d_i\in k$.
The element $\delta=1\otimes b-b\otimes1$
has $\delta^p=1\otimes a-a\otimes1=0$,
but is nonzero because its inverse image
$T-b$ has degree $1<p$ in the monic
quotient. Consequently the tensor of the
two original reduced fields is not reduced.
The same example shows that $L$ is not
geometrically reduced. Thus perfectness is
equivalent both to every reduced $k$-algebra
being geometrically reduced and to every
tensor product of two reduced $k$-algebras
being reduced.

For arbitrary original $k$-algebras $R,S$
over this perfect field, their nilradicals
$N_R=\operatorname{Nil}(R)$ and
$N_S=\operatorname{Nil}(S)$ give the exact
formula in $T=R\otimes_k S$
$$
\operatorname{Nil}(T)=J,
\qquad
J=\operatorname{im}(N_R\otimes_k S\to T)
+\operatorname{im}(R\otimes_k N_S\to T).
$$
This is a sum of ideals, with no assertion
that the sum is direct. The tensor quotient
map $T\to (R/N_R)\otimes_k(S/N_S)$ is
surjective and kills $J$. Its inverse
after passage to $T/J$ sends
$[r]\otimes[s]$ to $[r\otimes s]$.
Changing either representative changes
the tensor by an element of the indicated
ideal, so this is well-defined; balancing
and both composites follow on pure tensors.
Thus the kernel is exactly $J$.

Every element of $J$ is a finite sum
$z=z_1+\cdots+z_m$ of tensors having at
least one nilpotent factor. Choose positive
integers $d_i$ with $z_i^{d_i}=0$, and put
$D=1+\sum_{i=1}^m(d_i-1)$. Its full
multinomial expansion is
$$
z^D=
\sum_{\substack{a_1+\cdots+a_m=D\\a_i\geq0}}
\frac{D!}{a_1!\cdots a_m!}
z_1^{a_1}\cdots z_m^{a_m}=0.
$$
Indeed at least one $a_i\geq d_i$ in each
summand. The coefficients are the original
integers; no factorial is inverted in the
receiving ring. The empty sum is zero.
Hence $J\subseteq\operatorname{Nil}(T)$.
The proved quotient is reduced by the
first part, applied to $R/N_R$ and $S/N_S$.
Every nilpotent of $T$ therefore maps to
zero in that quotient and belongs to $J$,
proving the reverse inclusion and the
claimed formula with its exact quotient
isomorphism.
\end{proof}









